\section{Proof of the finite-detection limits}
\label{sec:fd-main}
We prove Theorem~\ref{thm:fd-limits}, including the estimates uniform over every zero and every allowed minor index.

\subsection{The analytic estimate that controls all heights}

We first obtain a relative convolution estimate uniform near every ray.
The theta integral gives, uniformly for $t$ in a compact set,
\[
 \log\max_{|z|\le R}|H_t(z)|=O(R\log(R+3)).
\]
For example, split its integral at $u=\log(R+3)$; beyond that point
$\exp(-\pi e^{4u})$ dominates the other exponential factors.
It follows that $F_a$ has order at most $1/2$.  Hadamard factorization
therefore has genus zero and has constant zero-free factor.  At $a=0$,
\[
 H_0(z)=H_0(0)\prod_{x>0}(1-z^2/x^2)^{m_x},\qquad
 N_H(R)=\sum_{0<x\le R}m_x\ll R\log(R+3).
\]
In this subsection all sums over $x$ include multiplicity.
For nonreal $z$ define
\[
 L(z)=\frac{H_0'(z)}{H_0(z)},\qquad
 S(z)=\sum_{x>0}\left(|z-x|^{-2}+|z+x|^{-2}\right).
\]
The denominator is nonzero by Alp\"oge's theorem.  The entire remainder
\[
 Q_z(v)=\frac{H_0(z+v)}{H_0(z)e^{L(z)v}}
\]
satisfies
\begin{equation}\label{eq:fd-rem}
 |Q_z(v)|\le e^{S(z)|v|^2/2},\qquad
 |Q_z(v)-1|\le8S(z)|v|^2e^{2S(z)|v|^2}.
\end{equation}
Indeed it is the locally uniform product of
$(1+v/(z-\xi))\exp(-v/(z-\xi))$, for $\xi=\pm x$.
The inequality $|1+w|^2\le\exp(2\Re w+|w|^2)$ proves the first bound.
For the second, scale so $S=1$.  The constant and linear Taylor
coefficients are $1,0$, and Cauchy's estimate on the unit circle gives
$|Q(v)-1|\le2e^{1/2}|v|^2$ on $|v|\le1/2$.
On its complement the first bound gives
$|Q(v)-1|\le2e^{|v|^2/2}\le8|v|^2e^{2|v|^2}$.

The heat equation and the Fourier integral give the exact Gaussian identity
\[
 H_{-a}(z)=\frac1{\sqrt{4\pi a}}
             \int_{\mathbb R}e^{-u^2/(4a)}H_0(z+u)\,du.
\]
Translate the line of integration by $p=2aL(z)$.  The integrand is entire,
and its Gaussian decay on the two vertical ends dominates the preceding
order-one growth bound.  Completion of the square yields
\[
 \frac{H_{-a}(z)}{H_0(z)e^{aL(z)^2}}=\mathbb E Q_z(p+U),
 \qquad U\sim N(0,2a).
\]
Set $A=aS(z)$ and $B=a^2|L(z)|^2S(z)$.  Equation~\eqref{eq:fd-rem},
$|p+U|^2\le2|p|^2+2U^2$, and elementary Gaussian integration give
\begin{equation}\label{eq:fd-gauss}
 \left|\frac{H_{-a}(z)}{H_0(z)e^{aL(z)^2}}-1\right|
 \le16e^{16B}\left(
 \frac{4B}{\sqrt{1-16A}}+
 \frac{2A}{(1-16A)^{3/2}}\right).
\end{equation}
This estimate applies for $16A<1$.  In particular $A,B\le2^{-10}$
make its right side at most $7/16$.  To check that last rational bound,
use $e^{16B}<2$, $(1-16A)^{-1/2}<2$, and
$(1-16A)^{-3/2}<3$.

For $r=|z|$ and $|\Im z|\ge qr$, $0<q\le1$, partial summation of
$N_H(R)\ll R\log(R+3)$ gives
\begin{equation}\label{eq:fd-cones}
 S(z)\ll q^{-2}\frac{\log(r+3)}{r+1},\qquad
 |L(z)|\ll q^{-1}\log(r+3).
\end{equation}
Here is the estimate without a separation hypothesis.  By symmetry take
$\Re z\ge0$ and split the positive zeros at $r/2$ and $2r$.
Below $r/2$ both distances are at least $r/2$, so their contribution
to $S$ is at most $8N_H(r/2)/r^2$.  The paired expression
$L(z)=-2z\sum_x(x^2-z^2)^{-1}$ gives a contribution
$O(N_H(r/2)/r)$ to $L$.  In the middle interval
$|z-x|\ge qr$ and $|z+x|\ge r$, giving respectively
$2q^{-2}N_H(2r)/r^2$ and $2q^{-1}N_H(2r)/r$.
For the last interval, $\sum_{x>2r}x^{-2}\ll\log(r+3)/r$,
which bounds the corresponding contributions by a constant times this
sum and by a constant times $r$ times this sum.  A fixed neighborhood
of zero is handled directly by the convergent product.

Define
\[
 \alpha(a)=\sup_{H_{-a}(z)=0,\,\Re z>0}|\arg z|.
\]
There are no imaginary-axis zeros, since $H_{-a}(iy)>0$.
The bounds~\eqref{eq:fd-cones} and~\eqref{eq:fd-gauss} show that
$\alpha(a)=O(\sqrt a)$.  More precisely,
\begin{equation}\label{eq:fd-tail}
 \lim_{R\to\infty}\sup_{0<a\le1}
 \sup_{\substack{H_{-a}(z)=0,\,\Re z>0\\|z|\ge R}}
            \frac{|\arg z|}{\sqrt a}=0.
\end{equation}
For completeness, fix $\eta\in(0,1/2)$ and put $q=\eta\sqrt a$.
In that cone the parameters in~\eqref{eq:fd-gauss} are bounded by
\[
 C\eta^{-2}\frac{\log(r+3)}{r+1},\qquad
 C^3\eta^{-4}\frac{\log^3(r+3)}{r+1}.
\]
Both tend to zero uniformly in $0<a\le1$.  Thus at sufficiently large
radius no zero has $|\Im z|\ge\eta\sqrt a\,|z|$.
The inequality $\arcsin u\le2u$ for $0\le u\le1/2$ proves
\eqref{eq:fd-tail}.  The same proof with a sufficiently large fixed
multiple of $\sqrt a$ gives the global $O(\sqrt a)$ bound.

At an endpoint of multiplicity $m$, the substitution
$a=q^2$, $z=x+qu$ gives, after division by its leading coefficient
and $q^m$,
\[
 e^{D_u^2}u^m+O(q)
   =(i\sqrt2)^m\operatorname{He}_m(u/(i\sqrt2))+O(q).
\]
All lowest weighted Taylor terms follow from
$\partial_a^jH_{-a}=\partial_z^{2j}H_{-a}$.
The Hermite roots are distinct: their Jacobi matrix has positive
successive off-diagonal entries $\sqrt1,\ldots,\sqrt{m-1}$, so an
eigenvector is determined by its first coordinate and each eigenspace
has dimension one.  The matrix is symmetric and therefore all roots
are real.  The implicit function theorem and Rouch\'e's theorem give
\begin{equation}\label{eq:fd-localheat}
 z_j(-a)=x+i\sqrt{2a}\,h_{m,j}+O(a).
\end{equation}
Simple endpoint zeros have real simple continuations by conjugation.
Compact zero convergence and~\eqref{eq:fd-tail} now prove
\begin{equation}\label{eq:fd-anglelimit}
 \lim_{a\downarrow0}\frac{\alpha(a)}{\sqrt a}
   =\sup_x\frac{\sqrt2\,h_{m_x}}x,
 \qquad h_m=\max_j|h_{m,j}|.
\end{equation}
The supremum is a maximum whenever it is positive.  In fact
$h_m\le2\sqrt m$, whereas $m_x\le N_H(2x)\ll x\log(x+3)$,
so its terms tend to zero at infinity.

\subsection{From angles to all Toeplitz minors}

We next use the elementary Schur identities to keep control of arbitrary indices.
For the alphabet $be^{i\theta},be^{-i\theta}$,
\[
 h_d=b^d\frac{\sin((d+1)\theta)}{\sin\theta},
\]
with value $(d+1)b^d$ at $\theta=0$.  The dual Jacobi--Trudi identity
turns a Toeplitz minor of order $r$ in elementary coefficients into
a skew Schur function with diagram width at most $r$.
One direct index formula, for allowed lists $i_\ell\le j_\ell$, is
\[
 \lambda'_p=j_{r+1-p}+p-i_1-r,\qquad
 \mu'_p=i_{r+1-p}+p-i_1-r\quad(1\le p\le r).
\]
Transposing the minor and reversing both lists gives
$s_{\lambda/\mu}$.  These formulas also prove the stated width bound.
A two-variable tableau has zero weight sum when a column has more
than two boxes.  Each two-box column otherwise contributes $b^2$,
and the unpaired interval in each row contributes one of the displayed
$h_d$, with $d\le r$.  The unpaired intervals are independent because
paired boxes below occupy an initial row segment and paired boxes above
occupy a final row segment.  Thus
\begin{equation}\label{eq:fd-sector}
 (r+1)|\theta|<\pi
 \quad\Longrightarrow\quad
 (1+be^{i\theta}w)(1+be^{-i\theta}w)\text{ has }PF_r\text{ coefficients}.
\end{equation}
Coefficient convolution is Toeplitz multiplication, so Cauchy--Binet
preserves this property under products.  For a fixed minor its sum
is finite, since intermediate indices stay between its extreme row
and column indices.  Local uniform limits preserve the inequalities.

The reciprocal parameters of $F_a$ are $z^{-2}$, one for each
right-half-plane zero of $H_{-a}$, counted with multiplicity.  Their
arguments have absolute values at most $2\alpha(a)$.
The genus-zero product and~\eqref{eq:fd-sector} therefore prove
$PF_r$ whenever $2(r+1)\alpha(a)<\pi$.
Every allowed minor is then strictly positive.  Indeed there are
infinitely many linear or conjugate quadratic factors, each with
constant coefficient one and positive first coefficient.  In iterated
Cauchy--Binet, move the rightmost row index to its target one step at
a time, then the next, using a fresh factor for every step.  The
selected factor minor is its positive first coefficient; all stationary
factor minors are one.  This gives a positive term, and all other terms
are nonnegative.  Finitely many factors realize the required path,
so the same positive lower bound survives the infinite limit.
There must be infinitely many factors because otherwise $H_{-a}$
would be a polynomial, whereas its Fourier integral tends to zero on
the real axis and has positive value at zero.

It follows already that $m_*(-a),R_*(-a)\gg a^{-1/2}$.
Equation~\eqref{eq:fd-anglelimit} proves that simplicity of all endpoint
zeros gives $\sqrt a\,m_*(-a)\to\infty$.
Edrei's entire-function theorem \cite[Proposition A, pp.~86--87;
Section~3, p.~92]{Edrei1953} identifies $PF_\infty$ with all zeros of
$F_a$ being negative real.  The map $w=-z^2$ identifies this with
all-real zeros of $H_{-a}$.  Alp\"oge and Rodgers--Tao give
$\Lambda=0$, so $m_*(-a)$ is finite for $a>0$.

\subsection{The orbital constant and isolation of one cluster}

The exact constant requires a whole coalescing cluster to be treated at once.
For a nonzero traceless Hermitian $D$ with sign-symmetric spectrum,
$\phi_D(X)=\int e^{i\operatorname{Tr}(DUXU^*)}\,dU$ is real and takes
a negative value.  On the $m^2$-dimensional Euclidean space of Hermitian
matrices with norm squared $\operatorname{Tr}(X^2)$, Gaussian Fourier
integration gives
\[
 \int e^{-t\|X\|^2/2}\phi_D(X)\,dX
   =(2\pi/t)^{m^2/2}e^{-\|D\|^2/(2t)}.
\]
The right side tends to zero as $t\downarrow0$.  Nonnegativity of
$\phi_D$ everywhere would contradict its being at least $1/2$ on a
fixed ball about zero.  A negative value can be diagonalized, perturbed
to distinct eigenvalues, and translated by a scalar without changing its sign.
This proves $\mathfrak r_m<\infty$.

There is strict positivity on every compact spectral-diameter range
$c<\mathfrak r_m$.  Differentiation under the integral gives
$\Delta\phi_m=-\|D_m\|^2\phi_m$.  At a zero in a nonnegative
neighborhood, the first nonzero homogeneous Taylor term would be
nonnegative and harmonic.  Its average on a sphere is its value at
zero, hence zero, contradicting its nonvanishing.  Such a neighborhood
therefore contains no zero.  Scalar translation puts every spectrum of
diameter at most $c$ in $[0,c]$, a compact set, giving a positive minimum.
The same observation at zero gives $\mathfrak r_m>0$.

We use the classical Harish--Chandra--Itzykson--Zuber integral
\cite[(11)--(12), pp.~312--313]{Zuber2018}:
\[
 \mathcal I(A,B)=\int_{U(m)}e^{\operatorname{Tr}(AUBU^*)}\,dU
  =\left(\prod_{j=1}^{m-1}j!\right)
       \frac{\det[e^{a_i b_j}]}{\Delta(a)\Delta(b)}.
\]
Repeated eigenvalues are interpreted by analytic continuation.
For $\rho=(m-1,\ldots,0)$ the Weyl alternant gives the exact identity
\begin{equation}\label{eq:fd-weylhciz}
 \frac{s_\mu(e^{\varepsilon a_1},\ldots,e^{\varepsilon a_m})}
      {s_\mu(1^m)}
 =\frac{\mathcal I(\operatorname{diag}(a),
                   \operatorname{diag}(\varepsilon(\mu+\rho)))}
        {\mathcal I(\operatorname{diag}(a),
                   \operatorname{diag}(\varepsilon\rho))}.
\end{equation}
This follows by dividing the two alternants; their remaining Vandermonde
quotient is $s_\mu(1^m)=\Delta(\mu+\rho)/\Delta(\rho)>0$.

Suppose a fixed cluster has parameters
$b_\varepsilon e^{\varepsilon a_j(\varepsilon)}$, with positive
$b_\varepsilon$ and
$a_j(\varepsilon)=i\kappa h_{m,j}+O(\varepsilon)$, and is stable
under conjugation.  For $c<\mathfrak r_m/\kappa$, formula
\eqref{eq:fd-weylhciz} is uniformly positive for every partition of
at most $m$ rows and width $r$ with $\varepsilon r\le c$.
The numerator ranges over a compact set of matrices with eigenvalues
in $[0,c+(m-1)\varepsilon]$ and tends uniformly to the strictly
positive orbital value.  The denominator tends to one.  Conjugation
makes the quotient real.  Partitions of more than $m$ rows have value
zero.  The Littlewood--Richardson expansion of a skew Schur function
is a sum with nonnegative integer coefficients of ordinary Schur
functions whose width is no larger than that of the skew diagram;
it therefore gives $PF_r$ for the whole cluster.

In particular $\mathfrak r_m\ge\pi/h_m$.  Apply
\eqref{eq:fd-sector} to the pairs in $e^{i\varepsilon h_{m,j}}$ and
take partitions with $\varepsilon\mu$ approaching an arbitrary
spectrum of diameter less than $\pi/h_m$.  Formula~\eqref{eq:fd-weylhciz}
shows that every orbital value in this range is nonnegative.
For $m=2$ the same formula is
$\phi_2(X)=\operatorname{sinc}(\lambda_1-\lambda_2)$.
For $m=3$, factoring its exponential Vandermonde gives
\[
 \phi_3(X)=\prod_{i<j}\operatorname{sinc}
              \left(\tfrac{\sqrt3}{2}(\lambda_i-\lambda_j)\right).
\]
The first formula gives $\mathfrak r_2=\pi$.  The second is positive
below diameter $2\pi/\sqrt3$, and negative at the spectrum
$(0,L/2,L)$ for $2\pi/\sqrt3<L<4\pi/\sqrt3$.
Thus $\mathfrak r_3=2\pi/\sqrt3$.

We now justify isolation of the negative cluster from every other zero.
Let its parameters be $\xi_j=b e^{\varepsilon a_j}$, with
$a_j=id_j+O(\varepsilon)$, $\sum a_j=0$, and distinct real $d_j$.
Let $\zeta_1,\ldots,\zeta_p$ be a finite conjugation-stable prefix
with positive limiting values strictly larger than $b^0=\lim b$.
Allow repeated prefix limits, but require pairwise separation at least
$c\varepsilon$ for distinct prefix parameters.  Let the remaining
parameters $\eta_j$ have uniformly bounded $\ell^1$ norm, modulus at
most $c_0<b^0$, and positive endpoint limits.  Their generating product
$C(v)=\prod_j(1-\eta_jv)^{-1}$ then converges locally uniformly on
$|v|<1/c_0$ to the positive endpoint product.

Choose distinct $L_1>\cdots>L_m>0$ at a strictly negative orbital
value of $\operatorname{diag}(d)$ and choose $C_1>L_1$.
Put
\[
 \mu_i=\lfloor L_i/\varepsilon\rfloor,\quad
 n=\lceil C_1/\varepsilon\rceil,\quad
 \lambda=(n^p,\mu_1,\ldots,\mu_m).
\]
Both $\mu_m$ and $n-\mu_1$ exceed a positive multiple of
$1/\varepsilon$.  The cluster character divided by
$b^{|\mu|}s_\mu(1^m)$ tends to the chosen negative orbital value,
by~\eqref{eq:fd-weylhciz}.

Here are the determinant estimates that transfer its sign.
Put $q=p+m$, $\mathcal A=\{\zeta_i,\xi_j\}$ and choose
$c_0<c_1<b^0$.  Partial fractions on $|v|=1/c_1$ give the full
complete symmetric coefficients in the form
\[
 h_k=\sum_{\alpha\in\mathcal A}A_\alpha\alpha^k+g_k,
 \quad |A_\alpha|\le C_2\varepsilon^{-(q-1)},\quad
 |g_k|\le C_2\varepsilon^{-(q-1)}c_1^k.
\]
Indeed the only small denominators are pairwise top-pole differences,
each bounded below by a constant times $\varepsilon$; after removing
the poles the same circle bounds the analytic remainder.
In the $q$ by $q$ Jacobi--Trudi determinant for $s_\lambda$, the
term using all top poles is exactly
\[
 s_\lambda(\mathcal A)\prod_{\alpha\in\mathcal A}C(1/\alpha).
\]
To check this equality, factor the pole matrix as
$[\alpha^{\lambda_i-i}]_{i,\alpha}
\operatorname{diag}(A_\alpha)[\alpha^j]_{\alpha,j}$.
The Vandermonde determinants cancel the internal factors in the residues.
The displayed product is real and has a positive nonzero limit.

Write $P=\prod_i\zeta_i>0$ and $W=P^n b^{|\mu|}>0$.
Every term using at least one remainder column uses distinct top-pole
labels for the other columns, since repeated labels give identical
columns.  Sort the logarithms of their moduli and assign them to the
nonincreasing row exponents.  Replacing any top modulus by $c_1$
loses a fixed factor less than one raised to at least $\mu_m$.
The chosen cluster moduli differ from $b$ by $e^{O(\varepsilon^2)}$,
and all prefix row exponents are $n$; hence the full top weight is
$W e^{O(\varepsilon)}$.
The fixed number of residue factors contributes only a power of
$\varepsilon^{-1}$.  Thus the whole remainder is
$O(W\varepsilon^{-K}q_1^{\mu_m})$, with $q_1<1$.

Expand the Weyl numerator of $s_\lambda(\mathcal A)$ along its
first $p$ rows.  The assignment of every prefix column to these
rows gives exactly
\[
 P^n s_\mu(\xi_1,\ldots,\xi_m)
       \prod_{i,j}(1-\xi_j/\zeta_i)^{-1}.
\]
The prefix block is $P^{n+m}\Delta(\zeta)$; cancellation against
the full Vandermonde proves the formula even with coalescing prefix
parameters.  Its extra product has a positive limit.
Every other assignment exchanges a larger prefix modulus and a
smaller cluster modulus between row lengths differing by at least
$n-\mu_1$.  Sorting again bounds the error by
$O(W\varepsilon^{-K'}q_2^{n-\mu_1})$, $q_2<1$.
Division by the Vandermonde costs only another power of
$\varepsilon^{-1}$.  Both errors are exponentially small in
$1/\varepsilon$, whereas the leading negative term has magnitude
at least a positive constant times $W s_\mu(1^m)$.
Consequently $s_\lambda<0$ for all sufficiently small $\varepsilon$.
Dual Jacobi--Trudi realizes it as a Toeplitz minor of order
$\lambda_1\le n$ in the window $0,\ldots,n+q-1$.
By translating the negative orbital spectrum, its smallest eigenvalue
and the extra gap $C_1-L_1$ can both be chosen arbitrarily small.
Thus the upper limit of $\varepsilon\lambda_1$ is at most the
negative spectral diameter plus an arbitrary positive error.

\begin{proof}[Proof of the Toeplitz limits in Theorem~\ref{thm:fd-limits}]
Put $\varepsilon=\sqrt a$.  Equation~\eqref{eq:fd-localheat} gives
near each endpoint $x$, after removing the positive geometric mean,
\[
 z_j(-a)^{-2}=b_x(a)e^{\varepsilon a_{x,j}},\qquad
 a_{x,j}=-\frac{2i\sqrt2}{x}h_{m_x,j}+O(\varepsilon).
\]
Let $\mathcal L=\min_{m_x\ge2}x\mathfrak r_{m_x}/(2\sqrt2)$.
Since $\mathfrak r_m\ge\pi/h_m$ and $h_{m_x}/x\to0$, these
local constants tend to infinity with $x$.  Every nonempty minimum
is therefore attained among finitely many endpoints.

Fix $c<\mathcal L$.  Choose a large circle so that~\eqref{eq:fd-tail}
makes every exterior parameter angle smaller than
$\pi\varepsilon/(2(c+1))$.  All exterior quadratic factors are
$PF_{\lfloor c/\varepsilon\rfloor}$ for small $\varepsilon$.
There are finitely many interior endpoint clusters, and the uniform
orbital argument gives the same property for each one; a simple
cluster is a positive linear factor.  Cauchy--Binet and the positive
path argument give strict positivity of every allowed minor of that
order.  This proves both lower limits, since $R_*\ge m_*-1$.
It also proves their divergence when $\mathcal L=\infty$.

For a multiple endpoint attaining $\mathcal L$, choose a negative
orbital spectrum within any prescribed positive error of its radius.
All preceding endpoint clusters form the finite prefix in the
isolation argument; Hermite simplicity gives their
$\varepsilon$ separation.  The complementary tail hypotheses hold
uniformly: the theta growth estimate and Jensen's formula give
$N_t(R)\ll R\log(R+3)$ on compact time intervals, and partial
summation makes $\sum_{|z_j(t)|>R}|z_j(t)|^{-2}$ tend uniformly
to zero.  Compact zero convergence handles the finite intervening
annulus and gives the required strict modulus separation.
Cluster isolation now produces negative minors of order $n$ in a
window of size $n+O_x(1)$, with
$\limsup\sqrt a\,n\le\mathcal L$ plus the arbitrary error.
This proves the upper limits and the first line of the theorem.
\end{proof}

\subsection{The exact Jensen degree}

We finish the proof by controlling both the finite endpoint clusters and the remaining zeros of every Jensen polynomial.
For an entire series $F(w)=\sum b_kw^k$ with $b_k\ge0$, its normalized
Jensen approximation obeys
\begin{equation}\label{eq:fd-japprox}
 \sup_{|w|\le R}|J_N(w/N)-F(w)|
 \le\frac{R^2F''(R)}{2N}\le\frac{F(2R)}{2N}.
\end{equation}
For $k\le N$, the coefficient multiplier is
$t_{N,k}=\prod_{j=0}^{k-1}(1-j/N)$ and
$0\le1-t_{N,k}\le k(k-1)/(2N)$.
For $k>N$ the same bound holds with $t_{N,k}=0$.
Sum the coefficient bounds and use $k(k-1)\le2^k$.
This proves~\eqref{eq:fd-japprox}.
Every $F_a$ with $a>0$ has a nonreal zero by the already established
Newman boundary.  Rouch\'e on a small disk around one such zero,
disjoint from the real axis, proves that $j_*(-a)<\infty$.

We need a finite-degree stability statement, for which the exact
input is Borcea--Br\"and\'en's Theorem~2(b), printed p.~470
\cite{BB2009}.  A real operator on polynomials of degree at most $N$
preserves hyperbolicity whenever its symbol $T[(z+w)^N]$ is nonzero
for $\Im z,\Im w>0$.
For $N\ge2$ and a conjugate parameter pair $be^{\pm i\theta}$, the symbol of
$1+2b\cos\theta D+b^2D^2$ is
\[
 (z+w)^{N-2}\big((z+w)^2+2Nb\cos\theta(z+w)+N(N-1)b^2\big).
\]
The quadratic discriminant is $4Nb^2(1-N\sin^2\theta)$.
Thus $N\sin^2\theta\le1$ and $|\theta|<\pi/2$ make its roots
negative real and its symbol stable.  A positive linear factor
$1+bD$ has stable symbol $(z+w)^{N-1}(z+w+Nb)$.
For $N=1$ the quadratic-factor operator has symbol
$z+w+2b\cos\theta$, so the same conclusion holds directly.
Their successive application to $z^N$ gives
\begin{equation}\label{eq:fd-operatorj}
 F(D)z^N=z^N J_N(1/z).
\end{equation}
Finite products preserve hyperbolicity, and coefficientwise limits of
monic hyperbolic polynomials remain hyperbolic.  The latter assertion
follows directly from Rouch\'e on a disk about any proposed nonreal
limit root.  We have proved that every parameter pair satisfying
$N\sin^2\theta\le1$ guarantees hyperbolicity of the degree-$N$
Jensen polynomial of the full product.

To examine a multiple endpoint, set
\[
 G(a,w)=H_{-a}(i\sqrt w)=\sum_{k\ge0}c_k(a)w^k,
 \qquad
 \mathcal J(a,\epsilon,w)=\sum_{k\ge0}c_k(a)w^k
                  \prod_{j=0}^{k-1}(1-j\epsilon).
\]
This interpolation is jointly entire.  Indeed the theta integral gives,
on $|a|\le A$,
$|G(a,w)|\le\exp(C_A(1+|w|^{3/4}))$.
Splitting the integral at $u=\log(\sqrt{|w|}+3)$ proves the stronger
bound $\exp(C_A(1+\sqrt{|w|})\log(\sqrt{|w|}+3))$.
Cauchy's estimate consequently gives
$|c_k(a)|\le C_A^k k^{-4k/3}$.
On $|\epsilon|\le E$ the multiplier product is at most
$(1+E)^k k!$, so the interpolating series converges normally on
compact sets.  In particular
\begin{equation}\label{eq:fd-jinterp}
 \mathcal J(a,1/N,w)=H_{-a}(0)J_{N,a}(w/N).
\end{equation}

The coefficient of $\epsilon^s$ in this interpolation is a differential
operator $Q_s$ of order at most $2s$, with principal part
\[
 \frac{(-1)^s w^{2s}}{2^s s!}D_w^{2s}.
\]
To verify it, let $e_s(k)$ be the elementary symmetric function of
$0,1,\ldots,k-1$.  Its recursion
$e_s(k+1)-e_s(k)=k e_{s-1}(k)$ shows inductively that it is a
polynomial of degree $2s$ with leading coefficient $1/(2^s s!)$.
Expand it in falling factorials of $k$; these are exactly the
coefficient multipliers of $w^jD_w^j$.
The heat equation becomes
\begin{equation}\label{eq:fd-wheat}
 \partial_aG=\mathcal L G,\qquad
 \mathcal L=-4wD_w^2-2D_w.
\end{equation}
Its $r$th Taylor coefficient is $\mathcal L^rG(0,w)/r!$ and has
principal differential term $(-4w)^rD_w^{2r}/r!$.

Suppose $G(0,w)=A(w-w_0)^m+O((w-w_0)^{m+1})$,
where $w_0=-x^2$ and $A\ne0$.
Substitute $a=\gamma q^2$, $\epsilon=q^2$, $w=w_0+qu$.
A term of bidegree $(r,s)$ in $a,\epsilon$ has differential order
at most $2(r+s)$, so no power of $q$ below $m$ occurs.
At degree $m$, only its principal differential coefficient contributes.
Summing over $r+s=k$ gives
\[
 \sum_{r+s=k}\frac{(-4\gamma w_0)^r}{r!}
                 \frac{(-w_0^2/2)^s}{s!}
       =\frac{(4\gamma x^2-x^4/2)^k}{k!}.
\]
Thus the quotient extends analytically to $q=0$ and
\begin{equation}\label{eq:fd-jlocal}
 \frac{\mathcal J(\gamma q^2,q^2,w_0+qu)}{Aq^m}
   =e^{c_x(\gamma)D_u^2}u^m+O(q),\qquad
 c_x(\gamma)=4\gamma x^2-x^4/2.
\end{equation}
The error is uniform on compact sets of $u,\gamma$ by joint analyticity.
When $m\ge2$ and $\gamma>x^2/8$, the leading polynomial has a
simple nonreal pair, since it is
$(i\sqrt{2c_x(\gamma)})^m
\operatorname{He}_m(u/(i\sqrt{2c_x(\gamma)}))$.
Rouch\'e preserves this pair.  Set $N=\lceil\gamma/a\rceil$,
$q=N^{-1/2}$, and replace $\gamma$ by $aN\to\gamma$ in
\eqref{eq:fd-jlocal}.  Equation~\eqref{eq:fd-jinterp} gives
\begin{equation}\label{eq:fd-jupper}
 \limsup_{a\downarrow0}a j_*(-a)\le x_*^2/8
 \quad\text{when }x_*<\infty.
\end{equation}

The matching lower bound needs more than compact Jensen approximation.
Fix $0<\gamma_0<x_*^2/8$ and choose
$\kappa=1/(4\sqrt{\gamma_0})$.
By~\eqref{eq:fd-tail}, choose $R_0$ so that all heat zeros of modulus
at least $R_0$ have $|\arg z|\le\kappa\sqrt a$ for $0<a\le1$.
Choose $R>R_0^2$ with no endpoint zero on $|w|=R$ and at least one
endpoint zero inside.  For small complex $a$ this contour contains
a fixed number $d$ of zeros of $F_a$.
The integrals
\[
 p_k(a)=\frac1{2\pi i}\int_{|w|=R}
             w^k\frac{\partial_wF_a(w)}{F_a(w)}\,dw
 \quad(1\le k\le d)
\]
are holomorphic and are their power sums.  Newton identities give a
monic polynomial with these zeros.  Divide by its nonzero constant
coefficient to obtain $P_a(0)=1$.
Then $F_a=P_a T_a$, with $P_a$ of degree $d$ and $T_a$ holomorphic
and nonvanishing inside the circle.  This factorization is jointly
holomorphic: inside, use the Cauchy integral of the boundary quotient;
outside, division is already holomorphic.  Both factors have real
coefficients for real $a$.

The finite polynomial $P_a$ has the same lowest weighted jets as $G$.
To see this, locally write $P_a=U(a,w)G(a,w)$, where $U$ is a
nonvanishing holomorphic unit.  At an endpoint $w_0$ of multiplicity
$m$, let $A_0$ be the leading coefficient of $P_0$ there.
For $2r\le m$, the coefficient of $a^r$ in $P_a$ is
\[
 \frac{A_0(-4w_0)^r m!}{r!(m-2r)!}(w-w_0)^{m-2r}
       +O((w-w_0)^{m-2r+1}).
\]
Indeed a term using a positive power $a^j$ from $U$ has vanishing
order at least $m-2r+2j$; the $j=0$ term has exactly the displayed
principal derivative from~\eqref{eq:fd-wheat}.  Variation of its
coefficient or of $U(0,w)$ adds at least one order.

Apply the finite interpolation
$\mathcal J_P(a,\epsilon,w)=\sum_{k=0}^d b_k^{P}(a)w^k
\prod_{j<k}(1-j\epsilon)$ to $P_a$.
The preceding weighted-jet argument proves~\eqref{eq:fd-jlocal}
for $\mathcal J_P$, with $A_0$ replacing $A$, uniformly for
$0\le\gamma\le\gamma_0$.
At a multiple endpoint, $x\ge x_*$ gives
$c_x(\gamma)<0$ uniformly on this compact interval; its limiting
polynomial is a real Hermite polynomial with simple separated roots.
At a simple endpoint the limiting polynomial is $u$.
Uniform Rouch\'e disks about these real roots each contain exactly
one root; real coefficients force that root to be real.  Their
$w$ coordinates remain negative.  The finitely many endpoint
multiplicities sum to $d$, so all roots of the degree-$d$ finite
interpolation have been counted.  Hence there is $N_1>d$ such that
\[
 N\ge N_1,\quad 0\le aN\le\gamma_0
 \quad\Longrightarrow\quad
 J_N(P_a)\text{ has only real negative zeros}.
\]

For the remaining finitely many $N<N_1$, strictness at $a=0$ is
needed.  Factor $P_0(w)=\prod_{i=1}^d(1+\beta_iw)$ with
$\beta_i>0$ and apply $1+\beta_iD$ successively to $z^N$.
The nonzero roots at each stage are negative and simple, with a
possible root at zero.  In fact, for a polynomial whose zero root
has multiplicity $v$ and whose other roots are simple negative
$\xi_1<\cdots<\xi_s$,
\[
 \frac{p'}p(z)=\frac vz+\sum_{j=1}^s\frac1{z-\xi_j}
\]
is strictly decreasing on each real complementary interval.
The equation $p'/p=-1/\beta$ has one simple solution to the left
of the smallest root and between consecutive distinct roots,
including the interval ending at zero when $v>0$.
No solution lies to the right of the largest root.  Old nonzero
roots cease to be roots; zero's multiplicity decreases by one.
This proves the induction and shows that every $J_N(P_0)$ has
only simple negative roots.  Coefficient continuity preserves them
for small real $a$.  Intersect these finitely many neighborhoods.
We have proved head hyperbolicity for every $N\le\gamma_0/a$.

Finally every conjugate parameter pair of $T_a$ has argument
$|\theta|\le2\kappa\sqrt a$ and therefore
\[
 N\sin^2\theta\le4\gamma_0\kappa^2=1/4
 \quad (N\le\gamma_0/a).
\]
Every real tail parameter is positive.  The already proved
finite-degree operator criterion shows that $T_a(D)$ preserves
hyperbolicity in degree at most $N$, first for finite tail products
and then by coefficientwise limits.  Apply it to the hyperbolic
polynomial $P_a(D)z^N$ and use~\eqref{eq:fd-operatorj}.
This gives hyperbolicity of $J_{N,a}$ for every $N\le\gamma_0/a$.
Thus
$\liminf_{a\downarrow0}a j_*(-a)\ge\gamma_0$.
Let $\gamma_0\uparrow x_*^2/8$ and combine with
\eqref{eq:fd-jupper}; for $x_*=\infty$ let $\gamma_0$ grow without
bound.  This proves the Jensen assertion and completes the proof
of Theorem~\ref{thm:fd-limits}.

\subsection{An explicit small-time certificate}

We give rational constants for the preceding lower bounds.
A direct ball evaluation of the entire completed function gives
\[
 0<p_0:=F_0'(0)=\frac{\xi''(1/2)}{8\xi(1/2)}<\frac12,
 \qquad
 p_0=0.00577624827885474269723345260758\ldots .
\]
The decimal is accompanied by an outward ball enclosure in the
\href{\workbenchurl}{\emph{Detection constants} certificate}.
The inequality, together with $\cos v\ge1-v^2/2$, proves
$H_0(x)\ge H_0(0)(1-p_0x^2)>0$ for $|x|\le1$.
Thus $x_1>1$ and $\sum_xx^{-2}=p_0<1/2$.

Hasanalizade--Shen--Wong's bound
\cite[Corollary~1.2, (1.5), p.~2]{HSW2021} is
\[
 \left|N_\zeta(T)-\frac{T}{2\pi}\log\frac{T}{2\pi e}\right|
 \le0.1038\log T+0.2573\log\log T+9.3675\quad(T\ge e).
\]
It gives $N_\zeta(e)<10$ and $N_\zeta(T)\le6T\log(T+3)$
for $T\ge e$.  Since $N_H(X)=N_\zeta(X/2)$, one obtains the
convenient bound $N_H(X)\le100X\log(X+3)$ for every $X>0$.
For $X<1$ the count is zero; for $1\le X<2e$ it is less than ten;
for $X\ge2e$ the preceding estimate is stronger.

For $r\ge1/2$, the three ranges in the proof of~\eqref{eq:fd-cones}
give the following respective upper bounds; write $\ell=\log(r+3)$:
\[
\begin{array}{c|cc}
 \text{range}&S&|L|\\ \hline
 x<r/2&400\ell/r&(400/3)\ell\\
 r/2\le x\le2r&800q^{-2}\ell/r&800q^{-1}\ell\\
 x>2r&3200\ell/r&(3200/3)\ell.
\end{array}
\]
For the last row, partial summation gives
$\sum_{x>2r}x^{-2}\le400\ell/r$.
Thus $S\le13200q^{-2}\ell/(r+1)$ and
$|L|\le2000q^{-1}\ell$.
For $r\le1/2$, the product and $p_0<1/2$ give
$S<4$ and $|L|\le4r/3$.
Both cases fit~\eqref{eq:fd-cones} with $C=2^{14}$.
The elementary inequalities
$\ell/(r+1)\le2$ and $\ell^3/(r+1)\le16$ give
$S\le2^{15}q^{-2}$ and $|L|^2S\le2^{46}q^{-4}$.
For the second inequality, use $r+1\ge(r+3)/3$ and the maximum
$\max_{y>0}(\log y)^3/y=27/e^3$ on $y\ge1$.
Taking $q=2^{16}\sqrt a\le1/2$ in~\eqref{eq:fd-gauss} proves
\[
 \alpha(a)\le2^{17}\sqrt a\qquad(0<a\le2^{-34}).
\]
Consequently, for every integer $r\ge1$, all allowed Toeplitz minors
of orders at most $r$ are strictly positive on
\[
 -\frac{2^{-36}}{(r+1)^2}\le t\le0.
\]
For $0<a\le2^{-38}$ this yields
\[
 m_*(-a)\ge\left\lfloor\frac{2^{-18}}{\sqrt a}\right\rfloor,
 \qquad
 R_*(-a)\ge\left\lfloor\frac{2^{-18}}{\sqrt a}\right\rfloor-1.
\]
The Jensen operator criterion gives, on the same time range,
\[
 j_*(-a)>\left\lfloor\frac{2^{-36}}a\right\rfloor.
\]
Indeed $|\theta|\le2\alpha(a)\le2^{18}\sqrt a\le1/2$,
so $N\sin^2\theta\le1$ throughout the displayed degree range.

\subsection{Computing the Hermite orbital radii}

The orbital constants can be reduced to finitely many one-variable equations.
Put
\[
 P_m(z)=i^m\operatorname{He}_m(z/i),\qquad
 S_m(x)=\sum_{j=1}^m\frac{e^{ih_{m,j}x}}{P_m'(ih_{m,j})},
 \qquad
 W_{m,k}(x)=\det[S_m^{(p+q)}(x)]_{p,q=0}^k.
\]
Then
\begin{equation}\label{eq:fd-wradius}
 \mathfrak r_m=
 \min_{0\le k\le\lfloor(m-2)/2\rfloor}
        \inf\{x>0:W_{m,k}(x)=0\}.
\end{equation}
Here a first-zero infimum is $+\infty$ when there is no positive zero.
To prove the formula, let $E_m$ be the real solution space of
$P_m(D)f=0$.  Lagrange interpolation gives
$S_m^{(j)}(0)=0$ for $j<m-1$ and $S_m^{(m-1)}(0)=1$.
Thus $u_j=S_m^{(m-1-j)}$, $0\le j\le m-1$, is a canonical basis
at zero: $u_j^{(i)}(0)=0$ for $i<j$ and $u_j^{(j)}(0)=1$.
The space is invariant under translations and reflections.
The confluent HCIZ formula says that $\phi_m$ vanishes at a spectrum
with distinct locations $t_\ell$ and multiplicities $\nu_\ell$ precisely
when the matrix with columns
\[
 \left((ih_{m,j})^r e^{ih_{m,j}t_\ell}/r!\right)_{j=1}^m,
 \qquad 0\le r<\nu_\ell,
\]
is singular.  This is the evaluation map from $E_m$ to the corresponding
Hermite jets.  Singularity is equivalent to a nonzero real solution
having at least $m$ zeros there, counted with multiplicity.
Consequently the critical interval length for the extended Chebyshev
property of $E_m$ is exactly $\mathfrak r_m$: below it all these
orbital values are strictly positive; beyond it a radial path from
zero to a negative value crosses a singular evaluation matrix.
At the endpoint take a convergent sequence of spectra with negative
values and decreasing diameters; its limit is a zero value, since a
negative limiting value could be radially contracted.

Carnicer--Mainar--Pe\~na's critical-length theorem
\cite[Section~3, Proposition~3, p.~6]{Carnicer2023} now applies.
It identifies the critical length of a reflection-invariant space with
the first positive zero among $W(u_j,\ldots,u_{m-1})$,
$j>(m-1)/2$.  Reversing its columns and putting $k=m-1-j$
gives exactly~\eqref{eq:fd-wradius}.
Also $\sin(h_mx)\in E_m$ has $m$ zeros on
$[0,(m-1)\pi/h_m]$, so
\[
 \frac\pi{h_m}\le\mathfrak r_m\le\frac{(m-1)\pi}{h_m}.
\]
For example $S_2=\sin x$ and
$S_3=(1-\cos(\sqrt3x))/3$.  The latter first zero is tangential,
which is why a sign-change search on the scalar functions alone
would be insufficient.

There is a rigorous algorithm for arbitrary rational accuracy in
\eqref{eq:fd-wradius}.  At zero,
\[
 W_{m,k}(x)=c_{m,k}x^{d}(1+O(x)),\quad
 d=(k+1)(m-1-k),\quad
 c_{m,k}=(-1)^{k(k+1)/2}\prod_{j=0}^k\frac{j!}{(m-1-j)!}.
\]
This follows by inserting the leading monomials in the determinant
and evaluating their Vandermonde.  Divide by $c_{m,k}x^d$.
The resulting entire function has value one at zero and computable
Taylor coefficients.  If
$W_{m,k}(x)=\sum_J A_Je^{b_Jx}$ is its Cauchy--Binet expansion,
then $J$ ranges over $(k+1)$-subsets,
\[
 A_J=\frac{\Delta((ih_{m,j})_{j\in J})^2}
             {\prod_{j\in J}P_m'(ih_{m,j})},\qquad
 b_J=\sum_{j\in J}ih_{m,j}.
\]
Writing $B=\sum_J|A_J|$ and $M=(k+1)h_m$, the degree-$N$
Taylor remainder of the normalized function on $|x|\le R$ is at most
\[
 \frac{BM^d}{|c_{m,k}|}
      \frac{(MR)^{N+1}e^{MR}}{(N+d+1)!},
\]
and its derivative is bounded by
$BM^{d+1}e^{MR}/|c_{m,k}|$.
A rational mesh therefore certifies positivity on every rational
$[0,c]$ with $c<\mathfrak r_m$.
Independently enumerate distinct rational spectra and certify a
negative orbital value by HCIZ.  Every rational
$C>\mathfrak r_m$ admits such a spectrum of diameter less than $C$,
by continuity.  Dovetail the two searches until their interval width
is smaller than the requested error.  Both searches terminate, including
when a scalar Wronskian has a tangential first zero.

Novak's uniform nonvanishing theorem supplies the dependence on $m$
\cite[Theorem~1.2, p.~6]{Novak2022}.  It gives an absolute
$\epsilon_0>0$ such that
$\int_{U(m)}\exp(zm\operatorname{Tr}(AUBU^*))\,dU\ne0$
for $|z|\le\epsilon_0$ and spectra of $A,B$ in the unit disk.
For a Hermitian $X$ of diameter $R$, take
$A=2(X-cI)/R$, $B=D_m/h_m$ and $z=iRh_m/(2m)$,
where $c$ is the spectral midpoint.  The integral is exactly $\phi_m(X)$.
It is real, starts at one on the path $uX$, and never vanishes when
$R\le2\epsilon_0m/h_m$.  Therefore
\[
 \frac{2\epsilon_0m}{h_m}\le\mathfrak r_m
                 \le\frac{(m-1)\pi}{h_m}.
\]
The Hermite Jacobi matrix gives
$\sqrt{m-1}\le h_m\le2\sqrt{m-1}$, the lower bound following
from its mean squared eigenvalue.
Its last $k$ coordinates have off-diagonal weights at least
$\sqrt{m-k}$.  The positive sine eigenvector of the unweighted
$k$-vertex path gives
$h_m\ge2\sqrt{m-k}\cos(\pi/(k+1))$.
Taking $k\to\infty$ with $k/m\to0$ proves $h_m/(2\sqrt m)\to1$.
In particular
\[
 \epsilon_0\sqrt m\le\mathfrak r_m\le\pi\sqrt{m-1},\qquad
 \epsilon_0\le\liminf\frac{\mathfrak r_m}{\sqrt m}
 \le\limsup\frac{\mathfrak r_m}{\sqrt m}\le\frac\pi2.
\]
The source theorem gives no numerical value of $\epsilon_0$ used here;
the rational small-time certificate above is independent of that constant.
